%! Mode:: "TeX:UTF-8"
%! TEX program = xelatex
\PassOptionsToPackage{quiet}{xeCJK}
\documentclass[withoutpreface,bwprint]{cumcmthesis}
\usepackage{tabularx}   % 解决 tabularx 环境缺失
\usepackage{float}      % 支持 [H] 选项
\usepackage{etoolbox}
\usepackage{ctex}
\BeforeBeginEnvironment{tabular}{\zihao{-5}}
\usepackage[numbers,sort&compress]{natbib}  % 文献管理宏包
\usepackage[framemethod=TikZ]{mdframed}  % 框架宏包
\usepackage{url}  % 网页链接宏包
\usepackage{subcaption}  % 子图宏包
\newcolumntype{C}{>{\centering\arraybackslash}X}
\newcolumntype{R}{>{\raggedleft\arraybackslash}X}
\newcolumntype{L}{>{\raggedright\arraybackslash}X}\renewcommand{\textbf}[1]{{\song\bfseries #1}}

\usepackage{tikz}
\usetikzlibrary{shapes.geometric, arrows.meta, positioning}

\tikzstyle{startstop} = [rectangle, rounded corners, minimum width=3cm, minimum height=1cm,text centered, draw=black, fill=red!20]
\tikzstyle{process} = [rectangle, minimum width=3.5cm, minimum height=1cm, text centered, draw=black, fill=blue!20]
\tikzstyle{decision} = [diamond, aspect=2, minimum width=3.5cm, minimum height=1cm, text centered, draw=black, fill=green!20]
\tikzstyle{arrow} = [thick,->,>=stealth]






















\title{基于集成学习的洪水概率预测}  % 论文标题



%%%%%%%%%%%%%%%%%%%%%%%%%%%%%%%%%%%%%%%%%%%%%%%%%%%%%%%%%%%%%
%% 正文
\begin{document}
	
% 封面和目录页不显示页码
\pagenumbering{gobble}
\maketitle
\begin{abstract}


\textbf{对于问题一,}

\textbf{对于问题二,}

\textbf{对于问题三,}

\textbf{对于问题四,}

最后，



\keywords{关键词；关键词；关键词；关键词；关键词}

\end{abstract}


%%%%%%%%%%%%%%%%%%%%%%%%%%%%%%%%%%%%%%%%%%%%%%%%%%%%%%%%%%%%% 
%\vspace{-1.4cm} 
%\tableofcontents  % 目录
%\newpage











%%%%%%%%%%%%%%%%%%%%%%%%%%%%%%%%%%%%%%%%%%%%%%%%%%%%%%%%%%%%% 
\pagenumbering{arabic}
\vspace{-1.4cm} 
\section{引言}
\subsection{问题背景}







%%%%%%%%%%%%%%%%%%%%%%%%%%%%%%%%%%%%%%%%%%%%%%%%%%%%%%%%%%%%% 
\subsection{研究意义}












%%%%%%%%%%%%%%%%%%%%%%%%%%%%%%%%%%%%%%%%%%%%%%%%%%%%%%%%%%%%% 
\subsection{问题重述}

\textbf{问题1：}  

\textbf{问题2：}  

\textbf{问题3：} 

\textbf{问题4：}  

















\newpage
%%%%%%%%%%%%%%%%%%%%%%%%%%%%%%%%%%%%%%%%%%%%%%%%%%%%%%%%%%%%% 
\section{总体分析}
























\newpage
%%%%%%%%%%%%%%%%%%%%%%%%%%%%%%%%%%%%%%%%%%%%%%%%%%%%%%%%%%%%% 
\section{模型假设}

为简化问题，本文做出以下假设：

\begin{itemize}[itemindent=2em]
\item 假设1
\item 假设2
\item 假设3
\end{itemize}

%%%%%%%%%%%%%%%%%%%%%%%%%%%%%%%%%%%%%%%%%%%%%%%%%%%%%%%%%%%%% 

\section{符号说明}









\begin{table}[H]
	\centering
	
	\begin{tabularx}{\textwidth}{CL}
		\toprule
		符号    & 说明    \\
		\midrule
		$n$ & 样本总数 \\
		$d$ & 每个样本的特征变量个数 \\
		$\boldsymbol{x}^{(i)}$ & 第 $i$ 个样本的 $d$ 个特征变量，$\boldsymbol{x}^{(i)} = (x_1^{(i)}, x_2^{(i)}, \dots, x_d^{(i)})$ \\
		$y^{(i)}$ & 第 $i$ 个样本对应的洪水发生概率 \\
		$f(\cdot)$ & 构建的预测函数 \\
		$\Omega(f)$ & 正则化项 \\

		\bottomrule
	\end{tabularx}
	
	\label{tab:符号说明}
\end{table}











\newpage
%%%%%%%%%%%%%%%%%%%%%%%%%%%%%%%%%%%%%%%%%%%%%%%%%%%%%%%%%%%%% 
\section{问题一的模型的建立和求解}
\subsection{具体分析}






\subsection{模型准备}





\subsection{模型建立}















\subsection{模型求解}
















\newpage
%%%%%%%%%%%%%%%%%%%%%%%%%%%%%%%%%%%%%%%%%%%%%%%%%%%%%%%%%%%%% 
\section{问题二的模型的建立和求解}
\subsection{具体分析}

\subsection{模型准备}

\subsection{模型建立}

\subsection{模型求解}




























\newpage
%%%%%%%%%%%%%%%%%%%%%%%%%%%%%%%%%%%%%%%%%%%%%%%%%%%%%%%%%%%%% 
\section{问题三的模型的建立和求解}
\subsection{具体分析}

\subsection{模型准备}

\subsection{模型建立}

\subsection{模型求解}
















\newpage
%%%%%%%%%%%%%%%%%%%%%%%%%%%%%%%%%%%%%%%%%%%%%%%%%%%%%%%%%%%%% 
\section{问题四的模型的建立和求解}
\subsection{具体分析}

\subsection{模型准备}

\subsection{模型建立}

\subsection{模型求解}











\newpage
%%%%%%%%%%%%%%%%%%%%%%%%%%%%%%%%%%%%%%%%%%%%%%%%%%%%%%%%%%%%%
\section{模型的分析与检验}

\subsection{灵敏度分析}

\subsection{误差分析}















\newpage
%%%%%%%%%%%%%%%%%%%%%%%%%%%%%%%%%%%%%%%%%%%%%%%%%%%%%%%%%%%%%
\section{模型的评价、改进与推广}

\subsection{模型的优点}
\begin{itemize}[itemindent=2em]
\item 优点1
\item 优点2
\item 优点3
\end{itemize}

\subsection{模型的缺点}
\begin{itemize}[itemindent=2em]
\item 缺点1
\item 缺点2
\end{itemize}




\subsection{模型的改进}






\subsection{模型的推广}






































\newpage
%%%%%%%%%%%%%%%%%%%%%%%%%%%%%%%%%%%%%%%%%%%%%%%%%%%%%%%%%%%%%
%% 参考文献
\nocite{*}
\bibliographystyle{gbt7714-numerical}  % 引用格式
\bibliography{ref.bib}  % bib源














\newpage
%%%%%%%%%%%%%%%%%%%%%%%%%%%%%%%%%%%%%%%%%%%%%%%%%%%%%%%%%%%%%
%% 附录
\begin{appendices}
	
	
	
	
\section{文件列表}
\begin{table}[H]
\centering
\begin{tabularx}{\textwidth}{LL}
\toprule
文件名   & 功能描述 \\
\midrule
%%%%%%%%%%%%%%%%%%%%%%%%%%%%%%%%%%%%%%%%%%%%%%%%%%%%%%%%%%%%%
数据预处理.py & 预处理程序代码 \\
q2.py & 问题二程序代码 \\
q3.c & 问题三程序代码 \\
q4.cpp & 问题四程序代码 \\
%%%%%%%%%%%%%%%%%%%%%%%%%%%%%%%%%%%%%%%%%%%%%%%%%%%%%%%%%%%%%
\bottomrule
\end{tabularx}
\label{tab:文件列表}
\end{table}








%\section{代码}
%\noindent 
%%%%%%%%%%%%%%%%%%%%%%%%%%%%%%%%%%%%%%%%%%%%
%数据预处理.py
%\lstinputlisting[language=python]{code/1.数据预处理.py}
%%%%%%%%%%%%%%%%%%%%%%%%%%%%%%%%%%%%%%%%%%%%%
%
%%%%%%%%%%%%%%%%%%%%%%%%%%%%%%%%%%%%%%%%%%%%
%q2.py
%\lstinputlisting[language=python]{code/q2.py}
%%%%%%%%%%%%%%%%%%%%%%%%%%%%%%%%%%%%%%%%%%%%
%
%
%q3.c
%\lstinputlisting[language=c]{code/q3.c}
%
%
%q4.cpp
%\lstinputlisting[language=c++]{code/q4.cpp}

\end{appendices}
\end{document}
